% Paper-ready methodology for the final frozen-classifier residual-GRU model.
% Required preamble packages:
% \usepackage{amsmath,amssymb,booktabs,array}
% Optional: \usepackage{microtype}

\section{Methodology}
\label{sec:methodology}

\subsection{Problem formulation}

We consider causal inference of hand state and reaching intent from four
surface electromyography (sEMG) channels and four six-axis inertial measurement
units (IMUs) distributed around the forearm.  At acquisition time $t$, the raw
wearable observation is
\begin{equation}
    \mathbf{x}_t = \left[\mathbf{e}_t,\mathbf{a}_t,\boldsymbol{\omega}_t\right],
\end{equation}
where $\mathbf{e}_t\in\mathbb{R}^{4}$ contains one sEMG measurement per sensor,
$\mathbf{a}_t\in\mathbb{R}^{12}$ contains three-axis acceleration from four
sensors, and $\boldsymbol{\omega}_t\in\mathbb{R}^{12}$ contains the corresponding
angular velocities.  The model produces, at 100 Hz, (i) a binary gripper state
$\hat{g}_t\in\{\text{open},\text{close}\}$, (ii) current Cartesian position
$\hat{\mathbf{p}}_t\in\mathbb{R}^{3}$, (iii) intended screen location
$\hat{\mathbf{u}}_t\in[0,1]^2$, (iv) a dense short-horizon trajectory
$\hat{\mathbf{p}}_{t+h}$ for $h=10,20,\ldots,200$ ms, and (v) a low-rate intent
trajectory at $h=100,200,\ldots,1000$ ms.  The long-horizon output is represented
as displacement from the current prediction,
\begin{equation}
    \hat{\mathbf{p}}^{\mathrm{intent}}_{t+h}
    = \hat{\mathbf{p}}_t + \Delta\hat{\mathbf{p}}_{t,h}.
\end{equation}
Only measurements at or before $t$ are available to the network.  Future VIVE
positions and future IMU samples are training targets and are never supplied as
inference inputs.

The final system is trained in two stages.  First, a strong causal
patch-transformer classifier is trained to recognize open and closed gripper
states.  Second, its modality encoders and classifier are frozen and reused by
lightweight state-aware residual-GRU adapters trained for current pose, screen
location, and future motion.  The raw signals are therefore encoded only once;
Stage~II adapts the shared features rather than instantiating another encoder.
Freezing the shared representation prevents motion-regression gradients from
degrading the gripper decision.

\subsection{Dataset organization and supervision}

Each trial is stored as a CSV file containing timestamps, four sEMG channels,
24 IMU channels, VIVE tracker position and orientation, a per-sample
\texttt{gripper\_state} field, and screen-target metadata when available.  The
gripper field is converted case-insensitively to class 0 (open) or class 1
(close); numerical encodings 0 and 1 are accepted equivalently.  Labels are
aligned to the uniform model grid using the most recent recorded label, subject
to a maximum age of 20 ms.  Samples without a sufficiently recent label are
masked rather than treated as either class.

VIVE measurements are used only as supervision.  Current position labels are
the synchronized tracker coordinates $\mathbf{p}_t$.  For the short-horizon
task, the target associated with time $t$ and horizon $h$ is
$\mathbf{p}_{t+h}$.  The long-horizon target is the displacement
$\mathbf{p}_{t+h}-\mathbf{p}_t$.  Screen targets are stored in normalized
coordinates $\mathbf{u}=(u_x,u_y)$, but conversion to pixels always uses the
horizontal and vertical dimensions independently:
\begin{equation}
    \mathbf{u}^{\mathrm{px}}
    = \left(Wu_x,Hu_y\right).
\end{equation}
This prevents the 16:9 display from being treated as a square coordinate
system; in our deployment $W=1920$ and $H=1080$.

Trial files are discovered recursively and exact byte-identical duplicates are
removed using SHA-256 hashes.  Unless a fixed external test participant is
specified, trials are shuffled with a fixed split seed and divided into 60\%
training, 20\% validation, and 20\% test subsets.  All windows from one trial
remain in the same subset.  Normalization statistics, class weights, early
stopping decisions, and operating thresholds are derived without using the test
set.  For the reported cross-participant application test, the new
participant's directories are withheld completely from both training stages
and evaluated only after the final checkpoint has been selected.  We
distinguish this external participant test from the internal trial-level test,
because a trial split alone is not evidence of participant-independent
generalization.

\subsection{Causal signal preprocessing}
\label{subsec:preprocessing}

All preprocessing is explicitly causal so that offline evaluation matches the
information available during streaming deployment.  Invalid timestamps are
removed; samples are stably sorted by acquisition time; and duplicate
timestamps retain only their final occurrence.  Missing sensor values are
filled from the most recent past observation for at most 20 ms.  Longer gaps
remain invalid and are represented by zeros accompanied by validity flags.

Each raw sEMG channel is filtered independently with a fourth-order causal
Butterworth band-pass filter.  The lower cutoff is 20 Hz and the upper cutoff
is
\begin{equation}
    f_{\mathrm{high}}=\min(450,0.4f_s)\;\mathrm{Hz},
\end{equation}
where $f_s$ is the acquisition rate.  Two trailing activation envelopes are
then formed using root-mean-square windows of 20 and 50 ms.  This produces eight
sEMG features---two scales for each of four physical channels.  No centered
filter, zero-phase filter, or future sample is used.

Each of the 24 IMU channels is first passed through a trailing three-sample
median filter and then a second-order causal 15 Hz Butterworth low-pass filter.
The DC component is retained so that gravity-related orientation information is
not discarded.  Both modalities are sampled onto a 100 Hz grid using the most
recent raw observation.  The resulting value dimensions are eight for sEMG and
24 for IMU.  Per-channel binary validity indicators are concatenated, giving
network inputs
\begin{equation}
    \tilde{\mathbf{e}}_t\in\mathbb{R}^{16},\qquad
    \tilde{\mathbf{i}}_t\in\mathbb{R}^{48}.
\end{equation}

Means and standard deviations are estimated channel-wise from valid training
samples only.  Standardized values are clipped to $[-20,20]$ before validity
flags are appended.  VIVE position is standardized independently per Cartesian
axis using training-only statistics, with a minimum standard deviation of
1 cm.  A time step participates in a fused task loss only when at least 75\% of
both its sEMG and IMU channels are recent and valid.

\subsection{Stage I: gripper-state classifier}
\label{subsec:classifier}

\subsubsection{Causal multiscale encoders}

The frozen classifier originates from the best validation checkpoint of the
neuromuscular-future model trained on the source participants.  It was first
optimized as part of the complete multitask network rather than trained as an
isolated linear classifier.  Consequently, the representation supporting the
open/close decision was simultaneously constrained by current pose, screen
target, endpoint, and future-motion supervision.  After checkpoint selection,
the entire classifier---not only its final linear layer---was frozen and reused
in Stage~II.

The classifier contains separate sEMG and IMU causal patch encoders; the two
raw modalities are not concatenated before temporal encoding.  This separation
is important because their physical meanings and useful temporal scales differ.
sEMG measures muscle recruitment and therefore supplies direct evidence about
whether the hand is being actively closed or opened.  IMU measurements describe
the resulting forearm motion and orientation and therefore provide mechanical
context.  For each modality, an input projection maps the packed features to
width $d=128$, followed by layer normalization and GELU activation.  A local
path applies causal depthwise convolutions with kernel sizes 5 and 9.  In
parallel, the projected sequence is divided into overlapping patches of length
16 frames and stride 4 frames.  Left padding ensures that the patch emitted at
time $t$ ends at $t$ rather than being centered on it.  Patch tokens receive
sinusoidal positional encodings and pass through four transformer layers with
four attention heads, feed-forward width $4d$, dropout 0.1, and a strictly
upper-triangular attention mask.  Patch features are held between patch
endpoints and combined with the frame-resolution local features.

Separate learned gates fuse sEMG and IMU at the context and local resolutions.
For modality features $\mathbf{h}^{e}_t$ and $\mathbf{h}^{i}_t$, a detached
gating input produces two softmax weights $\alpha^e_t$ and $\alpha^i_t$.  The
fusion layer receives not only the weighted features but also their difference
and product:
\begin{equation}
 \mathbf{z}_t = \phi\!\left(
  [\alpha^e_t\mathbf{h}^{e}_t,\,
   \alpha^i_t\mathbf{h}^{i}_t,\,
   \mathbf{h}^{e}_t-\mathbf{h}^{i}_t,\,
   \mathbf{h}^{e}_t\odot\mathbf{h}^{i}_t]
 \right).
\end{equation}
This supplies both modality-specific evidence and explicit cross-modal
agreement or disagreement.

\subsubsection{Why the state classifier uses both sEMG and IMU}

An sEMG-only classifier is appealing because opening and closing are generated
by muscle activation and, in a static-arm test, a fist contraction can often be
recognized without any inertial signal.  Static recognition, however, is not
the full deployment problem.  During reaching, transporting, and returning,
the same electrodes also measure wrist stabilization, forearm rotation,
co-contraction, and changes caused by electrode pressure or skin contact.  An
sEMG burst can therefore mean ``close the gripper,'' but it can also be a
movement-related activation that should not change the gripper state.  A model
that observes sEMG alone must infer this distinction entirely from a noisy and
participant-dependent activation pattern.

The IMU stream provides the missing context.  Acceleration and angular
velocity indicate whether the arm is stationary, accelerating through a reach,
rotating, dwelling near the object, or returning.  This information does not
replace sEMG: an IMU generally cannot observe an isometric fist contraction
when the forearm is stationary.  Instead, it answers a complementary question:
``under what mechanical condition did this muscle activation occur?''  The
classifier can then distinguish a contraction associated with grasp from a
similar-amplitude activation associated with transport or stabilization.

Our design therefore does not assume that either sensor is universally
dominant.  It implements three evidence paths:
\begin{enumerate}
    \item an sEMG encoder for direct neuromuscular evidence;
    \item an IMU encoder for movement phase, forearm dynamics, and orientation
          context; and
    \item an explicitly sEMG-only auxiliary alignment branch, described below,
          which prevents all state information from being delegated to the
          inertial stream.
\end{enumerate}
The main state head uses fused context and local features, whereas the auxiliary
branch must recover the state from sEMG.  Thus, fusion improves robustness in
motion without discarding the physiological origin of the gripper state.

The gates are evaluated independently at every time step and separately at the
local and contextual resolutions.  Local fusion can emphasize a short sEMG
transition when the fist changes state, while contextual fusion can retain IMU
evidence over a longer movement phase.  The features supplied to each gate are
detached from the modality encoders.  Gradients still train the gate itself,
but the encoders cannot artificially reshape their representations merely to
manipulate the gate.  The downstream fusion and task losses continue to train
both encoders normally.

This fused formulation should not be interpreted as proof that sEMG is useful.
A multimodal model can ignore one input even when that input is present.  We
therefore assess the physiological contribution by evaluating the same frozen
classifier with sEMG zeroed and with sEMG trials shuffled, in addition to
training dedicated sEMG-only and IMU-only baselines.  The static-arm clench test
demonstrates qualitative sensitivity to sEMG, whereas held-out degradation
under these interventions provides the quantitative evidence that sEMG adds
information beyond inertial context.

\subsubsection{Dense state head and auxiliary sEMG alignment}

The gripper classifier combines a stable context estimate with a
frame-resolution transition estimate.  Context logits are obtained from a
linear projection of the fused context.  Local logits are produced by a causal
depthwise convolution of kernel size 5 followed by pointwise convolutions.  The
combined logits are
\begin{equation}
    \boldsymbol{\ell}^{g}_t
    = \boldsymbol{\ell}^{\mathrm{context}}_t
    + \sigma(\beta)\boldsymbol{\ell}^{\mathrm{local}}_t,
\end{equation}
where $\beta$ is initialized to $-2$, initially favoring the stable context
path while allowing the local branch to sharpen transitions.

An additional causal sEMG--state alignment branch interleaves an sEMG token and
an action token at every time step.  During ordinary inference, the action
token is a learned MASK token, so the classifier cannot read the ground-truth
state.  The interleaved sequence is processed by a four-layer, four-head causal
transformer of width 64 with a bounded 100-frame attention context.  The branch
is trained to reconstruct masked sEMG activation and to classify masked state
tokens.  A zero-initialized projection adds its correction to the main gripper
logits, ensuring that the auxiliary branch is initially a no-op.

Although this checkpoint contains pose and future-motion heads, the final
system retains its gripper logits as the authoritative state estimate.  Thus,
the final state classifier is the completed fused sEMG+IMU classifier from
Stage I; it is not the unused private classifier inside the subsequent motion
GRU.

\subsubsection{Classifier training objective}

The principal classification term is class-balanced cross entropy,
\begin{equation}
 \mathcal{L}_{g}=-\frac{1}{|\mathcal{V}_g|}
 \sum_{t\in\mathcal{V}_g} w_{g_t}
 \log\operatorname{softmax}(\boldsymbol{\ell}^{g}_t)_{g_t},
 \qquad
 w_c=\frac{N}{2N_c},
\end{equation}
where $N_c$ is the number of valid training frames in class $c$.  The complete
Stage-I checkpoint was learned using multitask supervision:
\begin{align}
 \mathcal{L}_{\mathrm{cls-stage}}
 ={}& \mathcal{L}_{g}
 +0.2\mathcal{L}_{\mathrm{aux-state}}
 +0.15(\mathcal{L}_{\mathrm{masked-state}}+
        \mathcal{L}_{\mathrm{EMG-recon}})
 +0.03\mathcal{L}_{\mathrm{stable}} \nonumber\\
 &+0.2\mathcal{L}_{\mathrm{current-pos}}
 +0.5\mathcal{L}_{\mathrm{orientation}}
 +0.35\mathcal{L}_{\mathrm{pixel}}
 +0.2\mathcal{L}_{\mathrm{endpoint}}
 +0.1\mathcal{L}_{\mathrm{future-pose}}
 +0.05\mathcal{L}_{\mathrm{future-IMU}}.
\end{align}
The state-stability term penalizes consecutive changes in close probability
when adjacent labels are identical.  The future-pose decoder predicts every
10 ms through 200 ms, while the future-IMU objective asks the corresponding
latent to predict the change in IMU state.  These future quantities are targets
only.  This multitask training exposes the classifier to both muscle activation
and the mechanical context in which an open or closed state occurs.

Training uses AdamW with learning rate $3\times10^{-4}$, weight decay
$10^{-4}$, batch size 4, and gradient-norm clipping at 1.  The maximum training
length is 60 epochs with early stopping after 12 validation epochs without
improvement.  The selected checkpoint minimizes a validation score combining
current position error, orientation error, macro-F1, global and final-quarter
pixel errors, endpoint position, and endpoint orientation.  No test metric is
used for checkpoint selection.

\subsection{Stage II: shared-encoder neuromuscular residual-GRU adapters}
\label{subsec:residual-gru}

\subsubsection{Normalization bridge}

The frozen classifier and the motion heads may have been fit using different
training populations and therefore different normalization statistics.  It is
incorrect to feed motion-normalized values directly to the classifier.  Let
$z_m=(x-\mu_m)/\sigma_m$ denote the motion model input.  Before executing the
classifier, we analytically convert it to the classifier's standardized space:
\begin{equation}
 z_c = z_m\frac{\sigma_m}{\sigma_c}
       +\frac{\mu_m-\mu_c}{\sigma_c}.
\end{equation}
This affine bridge is applied separately to the eight sEMG features and the 24
IMU features.  Validity flags pass through unchanged.  The shared patch
encoders and classifier are kept in evaluation mode, all their parameters have
\texttt{requires\_grad=False}, and their outputs are computed without gradient
tracking.  Consequently, dropout cannot perturb the representation and
Stage-II losses cannot alter either the encoders or classifier.

\subsubsection{State-conditioned residual fusion}

Let $\mathbf{u}^{e}_t$ and $\mathbf{u}^{i}_t$ be the width-128 sEMG and IMU
context features from the frozen Stage-I patch encoders.  The motion network
applies separate two-layer unidirectional GRU adapters:
\begin{align}
 \mathbf{h}^{e}_{1:t} &= \mathbf{u}^{e}_{1:t}
       +W^e_a\operatorname{GRU}_e(\mathbf{u}^{e}_{1:t}),\\
 \mathbf{h}^{i}_{1:t} &= \mathbf{u}^{i}_{1:t}
       +W^i_a\operatorname{GRU}_i(\mathbf{u}^{i}_{1:t}).
\end{align}
Both adapters have hidden width 128 and dropout 0.1 between recurrent layers.
Their output projections $W^e_a$ and $W^i_a$ are initialized to zero, so the
initial adapters are identity mappings.  This combines the validated
multiscale classifier representation with recurrent motion dynamics without
duplicating the raw-signal encoders.
Let $\mathbf{q}_t=\operatorname{softmax}(\boldsymbol{\ell}^{g}_t)$ be the
detached probability from the frozen classifier and
$\mathbf{s}_t=\phi_s(\mathbf{q}_t)$ its learned embedding.  IMU is treated as
the mechanical base, whereas sEMG supplies a bounded state-dependent
correction:
\begin{align}
 \mathbf{c}_t &= \tanh\!\left(\phi_c([
      \mathbf{h}^{e}_t,\mathbf{s}_t])\right),\\
 \gamma_t &= \sigma\!\left(\phi_\gamma([
      \mathbf{h}^{i}_t,\mathbf{h}^{e}_t,\mathbf{s}_t])\right),\\
 \mathbf{f}_t &= \operatorname{LN}(\mathbf{h}^{i}_t+\gamma_t\mathbf{c}_t).
\end{align}
The final layer of $\phi_c$ is initialized to zero, and the gate's final layer
is initialized with zero weights and bias $-2$.  Training therefore begins at
the strong IMU representation rather than allowing noisy sEMG to overwrite it.
The network must learn an sEMG correction only where it reduces the supervised
task losses.  The correction is regularized to remain small.

\subsubsection{Current position and screen-intent heads}

Current Cartesian position is predicted by a two-layer MLP from
$\mathbf{f}_t$.  The screen head respects the observed multimodal structure of
the target space.  A classifier predicts one of nine cells in a $3\times3$
screen grid.  A second head predicts a bounded two-dimensional residual for
each cell.  With fixed normalized cell centers $\mathbf{m}_k$ and offsets
$\boldsymbol{\delta}_{t,k}\in[-1/6,1/6]^2$, the differentiable output is
\begin{equation}
 \hat{\mathbf{u}}_t = \sum_{k=1}^{9}
   \pi_{t,k}\operatorname{clip}_{[0,1]}
   (\mathbf{m}_k+\boldsymbol{\delta}_{t,k}),
 \qquad
 \boldsymbol{\pi}_t=\operatorname{softmax}(\boldsymbol{\ell}^{\mathrm{grid}}_t).
\end{equation}
Pixel errors are formed after scaling $x$ by $W$ and $y$ by $H$.  Later
observations receive greater weight through
$a_t=0.25+3.75r_t^4$, where $r_t\in[0,1]$ is progress through the recorded
trial.  This emphasizes convergence near the intended interaction without
discarding early-intent supervision.

\subsubsection{Short-horizon and one-second intent decoders}

For short-horizon trajectory prediction, each time step is paired with a
four-dimensional horizon encoding
\begin{equation}
 \mathbf{b}(\tau)=[\tau,\tau^2,\sin(\pi\tau),\cos(\pi\tau)],
\end{equation}
where $\tau=h/200\text{ ms}$.  A shared MLP maps
$[\mathbf{f}_t,\mathbf{b}(\tau)]$ to an absolute standardized Cartesian
prediction for each $h\in\{10,20,\ldots,200\}$ ms.

The one-second intent decoder is deliberately lower-rate.  At horizons
$h\in\{100,200,\ldots,1000\}$ ms, it receives fused mechanical state, the
explicit sEMG hidden state, and a horizon encoding:
\begin{equation}
 \Delta\hat{\mathbf{p}}_{t,h}
 = \phi_{\mathrm{intent}}([
      \mathbf{f}_t,\mathbf{h}^{e}_t,\mathbf{b}(h/1000\text{ ms})]).
\end{equation}
Including $\mathbf{h}^{e}_t$ explicitly prevents the long-horizon decoder from
accessing muscle information only through the small residual correction.  A
parallel training head predicts changes in low-rate IMU interval summaries
from $\mathbf{h}^{e}_t$ alone.  The current IMU feature is not connected
directly to this auxiliary decoder, reducing the shortcut of copying inertial
state and testing whether muscle activity predicts subsequent mechanics.
Raw future sEMG is not reconstructed because its high-frequency waveform is
poorly determined and would encourage the model to spend capacity on noise
rather than motion intent.

\subsection{Stage-II objective}

Let $\rho(\cdot)$ denote the element-wise smooth-$L_1$ loss and let
$\langle\cdot\rangle_{\mathcal{V}}$ denote a mean over valid elements only.  The
deployed task loss is
\begin{align}
 \mathcal{L}_{\mathrm{motion}} ={}&
  1.0\mathcal{L}_{\mathrm{pos}}
 +0.35\mathcal{L}_{\mathrm{pixel}}
 +0.15\mathcal{L}_{\mathrm{grid}}
 +0.10\mathcal{L}_{\mathrm{offset}} \nonumber\\
 &+0.50\mathcal{L}_{\mathrm{future200}}
 +0.10\mathcal{L}_{\mathrm{arrival-consistency}}
 +0.05\mathcal{L}_{\mathrm{intent-pos}}
 +0.05\mathcal{L}_{\mathrm{EMG\rightarrow future\text{-}IMU}} \nonumber\\
 &+0.01\mathcal{L}_{\mathrm{correction}}.
\end{align}
The authoritative classifier is frozen, so no Stage-II classification loss is
applied.  Likewise, orientation, trial-end endpoint, dense short-horizon IMU,
and long-horizon state losses are disabled in the final configuration.

The current-position loss is
\begin{equation}
 \mathcal{L}_{\mathrm{pos}}
 =\left\langle\rho(\hat{\mathbf{p}}_t-\mathbf{p}_t)\right\rangle.
\end{equation}
For the screen loss, normalized prediction errors are converted to units of
100 pixels independently on each axis before smooth-$L_1$ is evaluated and
weighted by $a_t$.  The grid term is nine-class cross entropy, while the offset
term supervises the residual associated with the ground-truth cell.

The 200 ms future loss averages valid smooth-$L_1$ errors over all 20 horizons:
\begin{equation}
 \mathcal{L}_{\mathrm{future200}}
 = \frac{1}{|\mathcal{H}_{200}|}
   \sum_{h\in\mathcal{H}_{200}}
   \left\langle\rho(\hat{\mathbf{p}}_{t+h}-\mathbf{p}_{t+h})\right\rangle,
 \quad \mathcal{H}_{200}=\{10,\ldots,200\}\text{ ms}.
\end{equation}
Prediction-to-arrival consistency additionally compares a future prediction
made at $t$ with the detached current-position prediction emitted when that
future time actually arrives:
\begin{equation}
 \mathcal{L}_{\mathrm{arrival-consistency}}
 = \frac{1}{|\mathcal{H}_{200}|}\sum_h
 \left\langle\rho(\hat{\mathbf{p}}_{t+h|t}
 -\operatorname{sg}(\hat{\mathbf{p}}_{t+h|t+h}))\right\rangle,
\end{equation}
where $\operatorname{sg}$ denotes stop-gradient.

For long-horizon intent, horizon contributions decay as
$\lambda_h=\exp(-h/500\text{ ms})$ and are normalized by
$\sum_h\lambda_h$:
\begin{equation}
 \mathcal{L}_{\mathrm{intent-pos}}
 = \frac{\sum_h\lambda_h
  \left\langle\rho(\Delta\hat{\mathbf{p}}_{t,h}
  -(\mathbf{p}_{t+h}-\mathbf{p}_t))\right\rangle}
 {\sum_h\lambda_h}.
\end{equation}
The interval-IMU target at horizon $h_k$ is the mean IMU state over the interval
$(h_{k-1},h_k]$ minus the IMU state at $t$.  This produces the analogous
$\mathcal{L}_{\mathrm{EMG\rightarrow future\text{-}IMU}}$ without requiring
reconstruction of every future high-rate sample.  Masked reconstruction of the
current sEMG waveform is disabled in the final shared-encoder model; the
cross-modal future-IMU target is the relevant representation objective.
Finally,
\begin{equation}
 \mathcal{L}_{\mathrm{correction}}
 = \left\langle\|\mathbf{c}_t\|_2^2\right\rangle
\end{equation}
discourages unnecessary sEMG corrections to the IMU base.

Stage II uses the same AdamW optimizer, learning rate, weight decay, gradient
clipping, batch size, maximum epoch count, and early-stopping patience as Stage
I.  Only the GRU adapters, state embedding, residual gate, and task decoders are
updated.  The best motion checkpoint minimizes the validation score
\begin{equation}
 S = E_{\mathrm{XYZ}} + 0.01E_{\mathrm{pixel}},
\end{equation}
because orientation and endpoint heads are disabled and frozen-classifier F1
is constant across Stage-II epochs.  Here $E_{\mathrm{XYZ}}$ is measured in cm
and $E_{\mathrm{pixel}}$ in pixels.

\subsection{Streaming inference}

At deployment, raw measurements are processed incrementally by the same causal
filters used during training.  The implementation retains the latest 2 s of
100 Hz features and requires a 300 ms warm-up.  At most one prediction is
emitted per 100 Hz frame.  If fewer than 75\% of either modality's channels are
recent, the output is marked invalid rather than silently fabricating a pose.
Standardized values are clipped exactly as during training, and diagnostics
report extreme normalization mismatch.

The frozen classifier supplies open and close probabilities.  A hysteretic
state machine closes the gripper when $P(\mathrm{close})\geq0.7$ and reopens it
when $P(\mathrm{close})\leq0.3$.  Between thresholds, the preceding state is
retained to prevent rapid switching.  Current Cartesian position is
unstandardized to metres.  The screen estimate is multiplied by
$(1920,1080)$.  The short-horizon decoder returns 20 absolute Cartesian points
through 200 ms; the intent decoder returns ten points through 1 s after adding
its predicted displacement to the current Cartesian estimate.

\subsection{Evaluation protocol}

Gripper performance is summarized by accuracy, macro-F1, and a fixed-order
open/close confusion matrix.  Macro-F1 is emphasized because it gives equal
importance to both states.  Current and future Cartesian errors are Euclidean
distances in centimetres.  Pixel performance is reported as Euclidean pixel
error together with absolute horizontal and vertical errors, since horizontal
error can dominate on a 1920-by-1080 display.  Median and 90th-percentile errors
are reported in addition to means.

Each future horizon is compared with a persistence baseline that holds the
model's current position prediction fixed:
\begin{equation}
 E_{\mathrm{hold}}(h)=
 \|\hat{\mathbf{p}}_t-\mathbf{p}_{t+h}\|_2,
 \qquad
 G(h)=E_{\mathrm{hold}}(h)-E_{\mathrm{model}}(h).
\end{equation}
A positive $G(h)$ therefore demonstrates predictive information beyond simply
reusing the current estimate.  Modality dependence should additionally be
tested by zeroing and trial-shuffling sEMG at inference, alongside separately
trained IMU-only and sEMG-only baselines.  Claims that sEMG contributes to
future intent are made only when these interventions worsen held-out
performance, not merely when an attention weight or residual gate is nonzero.

\subsection{Reproducibility configuration}

The final Stage-II model can be trained with the following command after the
Stage-I classifier checkpoint has been selected:
\begin{verbatim}
bash scripts/train_shared_encoder_residual_gru.sh \
  runs/gripper_neuromuscular_future_fused/emg_imu_best.pt \
  runs/shared_encoder_residual_gru_multi_person \
  data/shubham_open data/shubham_close \
  data/mukund_open data/mukund_closed \
  data/gazania_open data/gazania_closed
\end{verbatim}
The script fixes the model to fused sEMG+IMU input, two-layer causal GRU
adapters, 200 ms dense future supervision, 1 s intent reconstruction at 100 ms
intervals, EMG-to-future-IMU supervision, seed 42, a 100 Hz feature grid, and
the optimization settings described above.  The deployable checkpoint is
\texttt{emg\_imu\_best.pt}.  Repeated experiments should retain
the same split seed while varying the initialization seed, and should report
mean and standard deviation over at least three runs.
